Abstract


Public transportation systems in small and mid-sized cities often suffer from unreliable schedules, overcrowding, and limited operational visibility due to the absence of advanced and costly Intelligent Transport System (ITS) infrastructure. While large metropolitan regions deploy dedicated GPS devices and passenger-counting sensors, such solutions are financially impractical for resource-constrained transit agencies. This project presents a software-centric smart transit system designed to predict public transportation delays and improve service reliability by leveraging existing operational infrastructure such as Electronic Ticketing Machines (ETMs) and driver-side interfaces.

The proposed system integrates real-time bus tracking, occupancy estimation derived from ticketing transactions, and predictive analytics to generate accurate arrival time forecasts and crowd alerts. A centralized backend processes live ETM events, ticket data, and route information to support intelligent trip planning, dynamic fleet monitoring, and administrative decision-making. Passengers are provided with real-time bus locations, estimated arrival times, and occupancy indicators through a web-based interface, while administrators gain operational visibility through alerts and dashboards that enable proactive route and resource management.

The system architecture is supported by a fully normalized relational database designed up to Third Normal Form (3NF), ensuring data integrity, scalability, and efficient handling of both static and high-velocity data. By adopting a retrofit-first and zero-hardware approach, the project demonstrates that intelligent public transportation solutions can be implemented effectively without significant capital investment. The proposed model offers a scalable and cost-efficient pathway toward smarter, data-driven urban mobility, particularly suited for developing cities and emerging smart transportation ecosystems. order fulfillment, efficient inventory management, and a streamlined digital procurement process.












Table of Contents





												Page No.

	Acknowledgement									i

	Abstract                                                                                                                      ii
	
	Table of Content 									iii

	List of Figures									            iv
	
     1.	Introduction		                                                                                               1									
1.1	Objective………………………………………………………………………. 2	

1.2	Scope……………………………………………………………………………2

    2.    Software Requirement Specification………………………………………………. 3
	
	2.1 Software Requirements…………………………………………………………3
		
	2.2 Hardware Requirements………………………………………………………..3
	
	2.3 Functional Requirement………………………………………………………..5
	
	2.4 Non-Functional Requirements…………………………………………………5

   3.      Entity Relationship Diagram……………………………………………………… 6
	
   4.      Data Flow Diagram………………………………………………………………..11
   
	4.1 DFD Level 0…………………………………………………………………...12
	
	4.2 DFD Level 1……………………………………………………………………12

    5       Relational Schema and Normalization ……………………………………………14	
		
	1NF(First Normal Form)…………………………………………………………..15	
	
	2NF(Second Normal Form) ……………………………………………………….15
		
	3NF(Third Normal Form) …………………………………………………………15


  6.      Design of Forms,Security & Validation……………………………………………16
	
  7.      Conclusion………………………………………………………………………… 18	
	
           References………………………………………………………………………… 19

          Appendix: Snapshots……………………………………………………………… 21	

	
	








List of Figures

		
Figure No.	Figure Name	Page No.
1.		ER Diagram of Pharma Connect	8
2.		ER to Relational Model Mapping of Pharma Connect	9
3.		DFD Level 0 of Aetheria	12
4.		DFD Level 1 of Aetheria	13
5.		Login Page	21
6.		Admin Dashboard 	22
7.		Live Map	22
8.		Trip Planner 	23
9.		Driver Dashboard 	23
10.		ETM Controller 	24
11.		Admin Analytics 	25
12.		Results in User side	26
13.		Admin Analytics	27











 
CHAPTER 1

                                                Introduction
Public transportation forms the backbone of urban mobility; however, in many small to mid-sized cities, transit systems continue to suffer from unreliable schedules, overcrowding, and poor operational transparency. Unlike large metropolitan networks that can deploy costly GPS units, automated passenger counters, and proprietary fleet management systems, smaller municipalities operate under severe budget constraints and inconsistent ridership demand. This disparity often creates a vicious cycle—limited funding leads to inefficient operations, which in turn discourages ridership and further reduces revenue. 
To address these challenges, this project proposes a “Retrofit” Intelligent Transit System specifically designed for low-cost, low-infrastructure environments. Instead of relying on expensive, specialized hardware, the system leverages existing assets such as the driver’s smartphone and the Electronic Ticketing Machine (ETM) already deployed in buses. This retrofit approach enables modern transit intelligence without disrupting current operations or requiring major capital investment. 
The proposed system creates a seamless communication layer between passengers, drivers, and administrators. It provides real-time vehicle tracking, occupancy estimation derived from ticketing data, and dynamic resource allocation to adapt to unpredictable factors such as traffic congestion, variable dwell times, and fluctuating passenger demand. By transforming routinely collected operational data into actionable insights, the system enhances service reliability, improves passenger experience, and empowers transit authorities to make informed, data-driven decisions—ultimately fostering sustainable public transportation even in resource-constrained cities.









1.1	 Objective

•	Zero-Hardware Location Tracking: To design and implement a real-time vehicle location tracking mechanism that operates exclusively through the driver’s interface—such as a mobile application or Electronic Ticketing Machine (ETM)—thereby eliminating the need for dedicated external GPS hardware. 
•	Occupancy Awareness via Ticketing Data: To compute real-time bus crowding levels by analyzing ticketing transaction logs, including passenger boarding and alighting events, as a software-based alternative to physical passenger-counting sensors. 
•	Dynamic Fleet Reallocation: To enable transit administrators to identify high-demand and overcrowded routes in real time and dynamically reassign idle buses and drivers to these routes, improving service balance and overall fleet utilization. 
•	Smart Trip Planning: To provide passengers with an intelligent, multi-modal trip planning system that supports walking–bus–transfer combinations, while accounting for bus direction, route connectivity, and transfer feasibility to generate reliable end-to-end travel plans. 
•	Operational Visibility and Control: To deliver a comprehensive, real-time administrative dashboard that enables transport authorities to monitor fleet health, active and idle vehicles, route performance, and resource availability, thereby supporting data-driven operational decision-making. 

1.2	 Scope
The scope of this project extends significantly beyond basic real-time public bus tracking and serves as a foundation for a comprehensive, software-centric intelligent transportation ecosystem tailored for small to mid-sized cities. The system can be further enhanced by incorporating machine learning models to predict transit delays and arrival times more accurately using historical ticketing data, traffic patterns, weather conditions, and peak-hour demand trends. Advanced demand forecasting techniques can enable proactive fleet planning, allowing transport authorities to deploy buses in anticipation of congestion or special events rather than reacting to overcrowding after it occurs. The platform can also be scaled to support multi-modal transportation systems, including metro rail, feeder buses, and shared mobility services, thereby enabling seamless end-to-end trip planning for passengers. Additional scope includes real-time passenger notifications through mobile applications or SMS alerts, integration with city traffic signal systems for transit signal priority, and open API support for third-party mobility applications. From an administrative perspective, the system can evolve into a decision-support tool by offering long-term analytics, route performance benchmarking, and policy-level insights. Owing to its retrofit-based, zero-hardware architecture, the solution is highly adaptable, cost-effective, and replicable across diverse urban settings, making it a strong candidate for large-scale deployment in developing regions seeking sustainable, data-driven public transportation management..





CHAPTER 2

                          Software Requirement Specification

This chapter outlines the functional, non-functional, hardware, and software requirements necessary for the successful development and deployment of the Public Transit Management System. Clearly defining these requirements ensures a thorough understanding of the system’s objectives, scope, and operational expectations,addressing both the administrative needs of transit authorities and the information requirements of commuters.
The functional requirements describe the core features and services the system must provide, while the nonfunctional requirements specify quality attributes such as performance, security, reliability, and scalability.
Additionally, the hardware and software requirements identify the technological infrastructure needed for
implementation and operation. Together, these specifications establish a baseline for the system design,
development, and implementation phases, ensuring alignment with stakeholder expectations and real-world transit operations.


2.1 Software Requirements 

Frontend
o	Framework: React.js (Vite)
o	Styling: CSS
o	Scripting Language: JavaScript
Backend
o	Framework: Python (FastAPI)
o	Application Server: Uvicorn
Database
o	Relational Database: PostgreSQL
o	Deployment options: Neon DB (cloud) / Local PostgreSQL
o	Database Version: PostgreSQL 14
ORM (Object Relational Mapper)
o	Tool: SQLAlchemy
o	Used for Python-to-database interaction and query management
Map Integration
o	Libraries: Leaflet / React-Leaflet
o	Map Data Provider: OpenStreetMap
Development Tools
o	IDE / Code Editor: Visual Studio Code (VS Code)
o	Version Control: Git and GitHub
Package Management
o	Frontend: npm
o	Backend: pip (Python package manager)
Supported Browsers
o	Google Chrome (Latest)
o	Microsoft Edge (Latest)
o	Brave (Latest)



2.2 Hardware Requirements

The system is intentionally designed to operate with minimal hardware dependencies, following a retrofit-first approach suitable for resource-constrained environments. No dedicated GPS devices or passenger-counting sensors are required. On the server side, a standard machine with a minimum quad-core processor, 8 GB RAM, and 50 GB storage is sufficient to host the backend services, database, and simulation engine. On the client side, passengers and administrators require only a smartphone, tablet, or desktop device with internet connectivity and a modern web browser. Drivers interact with the system through an existing Electronic Ticketing Machine (ETM) or a smartphone, which provides ticketing input and periodic location updates. Optional hardware such as a mobile device GPS sensor may be utilized when available, but the system remains fully functional without it, ensuring low deployment and maintenance costs.


2.3 Functional Requirements


2.3.1 Passenger Module
•	The Passenger Module of the system is designed to provide commuters with real-time visibility into public bus operations. The system enables passengers to view live bus locations on an interactive map interface and observe current occupancy levels through intuitive color-coded indicators representing low, moderate, and high crowding conditions. It calculates and displays estimated arrival times at selected stops by continuously processing live and simulated movement data. Additionally, the system supports intelligent trip planning by identifying feasible routes, transfer points, and reverse-direction journeys, allowing passengers to plan reliable end-to-end trips based on real-time operational conditions..


2.3.2	Driver Module (ETM Interface)
2.3.3	The Driver Module, accessed through the Electronic Ticketing Machine (ETM) or a mobile interface, facilitates accurate data capture directly from on-ground operations. The system allows drivers to securely authenticate themselves and issue digital tickets to boarding passengers. Each ticketing transaction is time-stamped and processed by the backend to dynamically update bus occupancy levels. The driver interface also enables the submission of live location updates and operational status information. When manual driver input is received, the system automatically prioritizes this data over simulated values, ensuring that real-world events take precedence and maintaining operational accuracy.
2.3.4	Admin Module
The Admin Module provides transport authorities with centralized operational control and real-time situational awareness. The system presents administrators with a live dashboard displaying fleet status, including active, idle, and unavailable buses. It continuously monitors occupancy levels and passenger wait times to detect demand imbalances and generates alerts when predefined thresholds are exceeded. Administrators are empowered to dynamically reassign idle buses and drivers to overcrowded routes through the dashboard, with reassignment instructions instantly communicated to the respective driver interfaces. This functionality enables rapid response to disruptions and supports efficient fleet utilization.

2.3.5	Backend and Simulation Module
The Backend and Simulation Module serves as the core processing layer of the system, integrating data from passenger, driver, and administrative interfaces. The system simulates the concurrent movement of multiple buses to validate scalability and maintain continuous operation in the absence of manual updates. It employs mathematical interpolation to estimate smooth vehicle trajectories between stops and uses the Haversine formula to compute distances for nearest-stop detection and route planning. The backend ensures consistent synchronization between ticketing data, occupancy estimation, and real-time visualization across all user interfaces, enabling reliable and cohesive system behavior


2.4	 Non-functional Requirements

The system is designed to meet stringent non-functional requirements to ensure reliable and sustainable operation in real-world urban environments. Performance is a critical requirement, and the system is expected to support real-time tracking, ticketing updates, and dashboard visualization for multiple buses simultaneously with minimal latency. Scalability is emphasized to allow the seamless addition of new routes, vehicles, users, and cities without requiring fundamental architectural modifications. Reliability is ensured through consistent data synchronization mechanisms that prevent conflicts during concurrent ticketing, simulation, and administrative operations.

Security is addressed through role-based authentication and controlled access to system functionalities, ensuring that passengers, drivers, and administrators interact only with authorized features. Usability is prioritized by providing intuitive, lightweight interfaces suitable for users with varying levels of technical proficiency, particularly drivers operating ETMs in real-time conditions. Maintainability is achieved through modular backend services and structured APIs that simplify debugging, testing, and future enhancements. Portability is also a key requirement, enabling deployment across different operating systems and server environments, including local and cloud-based infrastructures, thereby supporting flexible and cost-effective adoption across diverse transit agencies.














CHAPTER 3

                                       Entity Relationship Diagram
An Entity-Relationship (ER) diagram is a conceptual representation used in Database Management Systems (DBMS) to model the structure of a database and the relationships among its entities. It provides a clear visual understanding of how data is organized and interconnected before implementation. ER diagrams help database designers and developers analyze data requirements, reduce redundancy, and ensure data integrity while designing the database schema.
In the Bus Kar Bhai project, the ER diagram is designed to accurately model the real-world operations of a modern public transport system involving passengers, drivers, buses, routes, stops, and digital ticketing. The objective of the ER model is to ensure data integrity, eliminate redundancy, and clearly define relationships between different entities. This structured approach helps in designing a normalized database schema that supports efficient querying, real-time location tracking, and reliable high-volume transaction processing.
3.1	 Entities used:
The following entities are used to organize and manage data in the Moovit-Chalo system:
•	User Entity: Stores details of all system actors, including Passengers, Drivers, and Administrators. It maintains attributes such as user ID, name, email, password hash, role (User/Driver/Admin), and current driver status.
•	Bus Entity: Maintains details of the physical fleet, including bus number, total capacity, current condition, active status, and manual override status. It represents the core asset tracked by the system.
•	Route Entity: Contains information about specific transit paths, including route ID, route code (e.g., 500-A), source location, and destination location. It defines the abstract path a bus travels.
•	Stop Entity: Stores geospatial data for transit points, including stop ID, stop name, latitude, and longitude. It represents the physical locations where passengers board or deboard.
•	Ticket Entity: Records transaction details for passengers, including ticket ID, source stop, destination stop, passenger count, ticket timestamp, and the associated bus ID.
•	Schedule Entity: Manages the planned timings for routes, including departure times and day types (Weekday/Weekend), linking specific routes to time slots.
•	Notification Entity: Manages alerts sent to users, such as "Bus Arriving" or "Route Assigned," including message content, read status, and creation timestamps.
These entities collectively support the core functionality of fleet management, live tracking, dynamic routing, and digital ticketing within the Bus Kar Bhai system.







3.2	 Relationships uses:
Relationships define how entities interact within the system:
•	A User (Passenger) can purchase multiple Tickets (1:N relationship), while each ticket is associated with a single user account.
•	A Route consists of multiple Stops, and a Stop can be part of multiple Routes, forming a Many-to-Many (M:N) relationship. This is managed through an associative entity (Route_Stops) that also stores the sequence number and distance.
•	A Bus is assigned to one active Route at a time (N:1 relationship), but a route can have multiple buses running on it simultaneously.
•	A Bus can generate multiple ETM Events (1:N relationship), such as GPS updates or occupancy changes, which are logged for historical tracking.
•	A User (Driver) is associated with a Bus during a shift, often represented dynamically through status updates or assignments.
•	An Admin can send multiple Notifications (1:N relationship), while each notification belongs to a specific target user.
These relationships ensure efficient data organization, consistent real-time tracking, and reliable coordination between the fleet, the passengers, and the administrative control center.
3.3	 ER Diagram: 
The ER diagram serves as a fundamental design tool for the Moovit-Chalo database system. It visually represents key entities such as users, buses, routes, stops, and tickets, along with the relationships among them. By separating static data (like Routes and Stops) from dynamic data (like Live GPS Events), the design minimizes data redundancy and improves query performance for real-time maps.
The use of an associative entity for Route-Stops ensures accurate handling of complex transit networks where multiple routes share common stops. Primary and foreign keys enforce referential integrity, while normalization up to Third Normal Form (3NF) ensures efficient storage and prevents data anomalies. The ER model also supports future scalability, allowing features such as predictive analytics, pass systems, and multi-modal transport to be integrated with minimal schema changes, providing a robust foundation for reliable system implementation.
.
 


Figure 1: ER Diagram 
3.4	 Relational Mapping from ER Diagram
Relational mapping is the process of converting the conceptual Entity-Relationship (ER) diagram into a set of well-defined relational tables that can be implemented in a relational database system. In the Moovit-Chalo project, each entity in the ER diagram is mapped to a corresponding table with a primary key to uniquely identify records. Relationships such as one-to-many and many-to-many are implemented using foreign keys, ensuring referential integrity between tables.
In this system, core entities like Users, Buses, Routes, and Stops are transformed into normalized tables. One-to-many relationships, such as a User purchasing multiple Tickets or an Admin sending multiple Notifications, are handled by embedding the User ID as a foreign key in the child tables.
Associative entities like Route_Stops are specifically designed to resolve the complex many-to-many relationship between Routes and Stops, allowing a single stop to be part of multiple routes with specific sequence numbers. Additionally, the ETM_Events table links high-velocity dynamic data (speed, occupancy) to specific Buses and Routes using foreign keys. This mapping ensures a normalized database structure that minimizes redundancy, supports efficient querying for live location updates, and provides a reliable foundation for integrating static schedule data with real-time fleet operations.
                       

Figure 2: ER to Relational Model mapping 








CHAPTER 4

Data Flow Diagram
A Data Flow Diagram (DFD) is a graphical representation used to model the flow of data within a system and to illustrate how information is processed, stored, and exchanged between different system components. It provides a clear and structured way to visualize system functionality without focusing on implementation details such as programming logic or hardware configuration. DFDs help in understanding system requirements, identifying data dependencies, and analysing the interaction between external entities, processes, and data stores.
In this project, Data Flow Diagrams are used to represent the operational behaviour of the Predict Public Transportation Delay System at multiple levels of abstraction. The diagrams illustrate how passenger inputs, sensor data, and administrative commands flow through the system to enable real-time tracking, delay prediction, and intelligent trip planning. The hierarchical decomposition from Level-0 (context diagram) to Level-1 and Level-2 provides increasing detail, making it easier to understand both the overall system boundary and the internal processing logic. This structured representation ensures clarity, supports effective system analysis, and facilitates better communication between developers, stakeholders, and evaluators.

4.1 DFD Level 0
The Level-0 Data Flow Diagram represents the system as a single unified process, labeled as the BMTC AI Smart Transit System, and illustrates its interaction with external entities. The primary external entities involved are the Passenger, Admin, and Bus Sensors (GPS/IoT). Passengers provide inputs such as destination details and current location to the system, and in return receive trip plans, estimated time of arrival (ETA), and real-time alerts related to delays and crowding. The Bus Sensors supply live latitude, longitude, and occupancy-related data, which form the core input for tracking and prediction. The Admin interacts with the system by issuing dispatch commands and receives crowd alerts, performance statistics, and operational insights. This level establishes the system boundary and highlights how the intelligent transit system acts as a centralized platform that integrates user requests, sensor data, and administrative control to deliver real-time transit intelligence.
.
 
Figure 3: DFD Level 0 
4.2	DFD Level 1
The Level-1 Data Flow Diagram decomposes the overall system into major functional processes and associated data stores, providing a clearer view of internal data movement. The Trip Planner (Process 1.0) handles passenger search queries and uses static route and stop information stored in D1 (Routes/Stops) to generate feasible travel plans. The Live Tracking Engine (Process 2.0) continuously receives raw GPS data from bus sensors and updates real-time positional information in D2 (Live State / Coordinates). This live data is forwarded to the AI Prediction Module (Process 3.0), which analyzes movement patterns and delay trends to generate ETA predictions and alerts stored in D3 (Predictions). The Admin Dispatcher (Process 4.0) reads predictive alerts and system status to enable critical decisions such as bus dispatching and resource reallocation, while also allowing administrators to issue control commands back into the system. Together, these processes demonstrate how tracking, prediction, planning, and administration are tightly integrated to ensure efficient and responsive transit operations.

 
Figure 4: DFD Level 1 

4.3	 DFD Level-2: Trip Planner Logic (Process 1.0) Description

                The Level-2 Data Flow Diagram provides a detailed breakdown of the internal logic of the Trip Planner module. The process begins with Resolve Locations (1.1), where user-provided input is converted into precise latitude and longitude coordinates. These coordinates are passed to Filter Routes (1.2), which accesses the D1 Routes/Stops data store to identify relevant bus routes that serve the requested locations. The resulting route identifiers are forwarded to Fetch Prediction (1.3), which retrieves estimated arrival times and delay information from the D3 Predictions data store. The Rank Options (1.4) process then evaluates and prioritizes available travel options based on ETA, predicted delays, and user preferences stored in D5 User Favorites. Finally, Generate Response (1.5) compiles the ranked results into a structured trip plan that is displayed to the passenger. This level clearly illustrates how static route data, live predictions, and personalized preferences are combined to deliver intelligent, reliable trip planning.
                                  

Figure 5: DFD Level 1 


CHAPTER 5	

Relational Schema and Normalization
In the relational model of a Database Management System (DBMS), a schema defines the logical structure of the database by specifying tables, their attributes, and the relationships between them. A relational schema describes how data is organized and stored in the form of relations, where each table consists of rows (tuples) and columns (attributes). Columns represent specific data fields, while rows represent individual records. This structured representation ensures efficient storage, retrieval, and management of data while maintaining consistency and integrity across the system.
In the Moovit-Chalo platform, the relational schema is designed to support a centralized public transport management ecosystem. The schema organizes data related to passengers, bus drivers, vehicle fleets, routes, stops, schedules, and live tracking events, enabling efficient coordination between multiple stakeholders. This structured design ensures accurate ticket transaction processing, reliable real-time tracking, and seamless integration between relational data storage and cloud-based microservices, while maintaining data consistency and integrity across the entire transit workflow.
.
5.1 Relational Model
          The relational model for Moovit-Chalo is derived by mapping the Entity-Relationship (ER) diagram to relational tables. Each entity in the ER diagram is converted into a corresponding table, with attributes represented as columns and a primary key assigned to uniquely identify each record. Relationships between entities are implemented using foreign keys based on their cardinality.
One-to-many relationships are represented by including the primary key of the parent entity as a foreign key in the related child entity. For example, the relationship between User and Ticket is implemented by storing user_id as a foreign key in the tickets table, and the relationship between Route and Schedule is represented by including route_id as a foreign key in the schedules table. Many-to-many relationships, such as between Route and Stop, are resolved using an associative table named route_stops, which includes foreign keys referencing both route_id and stop_id along with sequence attributes. Dependent entities like Notification rely on foreign keys (user_id) to maintain referential integrity with their parent entities. This relational mapping ensures a well-structured, consistent, and normalized database design that accurately reflects real-world public transport operations and is suitable for implementation in a PostgreSQL database system
.
.
5.2 Normalization
Normalization is the process of organizing database tables to reduce data redundancy and improve data integrity. Redundancy occurs when the same data is stored in multiple places, leading to inconsistencies and update anomalies. Normalization addresses these issues by decomposing tables into smaller relations and organizing them according to a series of normal forms.

1.	5.2.1 First Normal Form (1NF)
A relation is said to be in First Normal Form (1NF) if all its attributes contain atomic and indivisible values, and there are no repeating groups or multi-valued attributes. In the Bus Kar Bhai database schema, all tables store single-valued attributes such as user details, bus specifications, route codes, and stop locations. Complex data types like geospatial coordinates are split into atomic lat and lon columns, and lists of stops are handled via separate rows in the route_stops table rather than comma-separated lists. Each relation is defined with a clearly identified primary key, and no attribute contains multiple or composite values. Hence, all relations in the Bus Kar Bhai system satisfy the conditions of First Normal Form (1NF).

2.	5.2.2 Second Normal Form (2NF)
A relation is in Second Normal Form (2NF) if it is already in First Normal Form and does not contain any partial dependency on a composite primary key. In the Bus Kar Bhai schema, most relations have single-attribute primary keys (e.g., id or bus_number), eliminating the possibility of partial dependency. In relations with potential composite keys, such as route_stops (which effectively relies on the combination of route_id and stop_id), all non-key attributes like sequence_number and distance_km depend on the complete composite key pair and not on a subset of it. Therefore, the database schema satisfies the conditions of Second Normal Form (2NF).

3.	5.2.3 Third Normal Form (3NF)
The Third Normal Form (3NF) in the Bus Kar Bhai database schema was achieved by identifying and eliminating transitive dependencies to ensure that all non-key attributes depend only on the primary key. In the initial design, attributes such as route_source and route_destination might have been transitively dependent on schedule_id if stored there. By separating ROUTES and SCHEDULES into independent relations, all route-related attributes were made directly dependent on route_id, eliminating indirect dependencies.
Similarly, in the TICKETS relation, attributes such as ticket_timestamp and passenger_count depend solely on the Ticket id. Bus details like capacity are referenced via bus_id rather than being stored redundantly in the ticket record. In the BUSES and ETM_EVENTS relations, static attributes (Capacity) and dynamic attributes (Speed, Current Occupancy) are separated. This ensures that high-frequency updates to live tracking data do not require duplicating static vehicle information. Pricing attributes in the ROUTES table are directly dependent on the primary key, ensuring no non-key attribute depends on another non-key attribute.

The final normalized relational schema for the Bus Kar Bhai system complies with First Normal Form (1NF), Second Normal Form (2NF), and Third Normal Form (3NF). This normalized design minimizes data redundancy, improves data integrity, and provides a scalable and reliable foundation for efficient database operations in a real-world public transport management system.
.



CHAPTER 6

Design of Forms, Security and Validation

T In the Bus Kar Bhai cloud-based transport system, data integrity and user security are paramount. This section details the interface design for data collection, the security protocols implemented to protect sensitive information, and the validation logic used to ensure data accuracy.

6.1 Design of Forms
Forms serve as the primary interface for users to interact with the system. The application uses React.js to create dynamic, responsive forms that adapt based on the user role (Passenger, Driver, or Admin).
Key Forms in the System:
1.	User Authentication Form (Login/Register):
o	Fields: Full Name, Email Address, Password, Role Selection (User/Driver).
o	Design: A toggle-based unified card design (as seen in the Login module) that switches between "Sign In" and "Sign Up" modes to reduce UI clutter.
o	Purpose: Captures credentials to generate a secure session token.
2.	Trip Planning & Search Form:
o	Fields: Source Location (Dropdown/Search), Destination Location (Dropdown/Search).
o	Design: Located prominently on the passenger dashboard. It uses auto-suggest dropdowns populated from the Stops database to prevent invalid location entries.
3.	Ticket Booking Form:
o	Fields: Bus Selection (Auto-filled), Source, Destination, Passenger Count.
o	Design: A clean modal or card interface on the Passenger and Driver apps. It dynamically calculates the fare based on the selected route distance.
4.	Admin Fleet Allocation Form:
o	Fields: Select Route, Select Available Bus, Select Driver.
o	Design: A dashboard control panel that filters only IDLE buses and drivers, preventing the assignment of already busy resources.
________________________________________
6.2 Security Mechanisms
The Bus Kar Bhai system implements a multi-layered security architecture to protect user data and preventing unauthorized access.
1.	Authentication via JSON Web Tokens (JWT):
o	Instead of traditional server-side sessions, the system uses stateless JWTs.
o	Process: Upon successful login, the FastAPI backend issues a signed token containing the user's ID and Role. This token is stored in the browser's localStorage and sent in the Authorization header (Bearer Token) for every subsequent API request.
o	Benefit: Allows the backend to be stateless and easily scalable across multiple cloud containers.
2.	Password Hashing (Bcrypt):
o	Implementation: User passwords are never stored in plain text. The system uses the Bcrypt hashing algorithm (via the passlib library) to encrypt passwords before saving them to the PostgreSQL database.
o	Verification: During login, the entered password is hashed and compared against the stored hash, ensuring that even if the database is compromised, actual passwords remain hidden.
3.	Cross-Origin Resource Sharing (CORS):
o	Since the Frontend (React) and Backend (FastAPI) run on different ports/domains, strict CORS policies are configured.
o	Policy: The backend is configured to accept requests only from the trusted frontend origin (e.g., http://localhost:5173 or the deployed cloud URL), blocking malicious requests from unauthorized domains.
4.	SQL Injection Prevention:
o	The system uses SQLAlchemy ORM (Object-Relational Mapping) for all database interactions.
o	Mechanism: Instead of raw SQL strings, queries are constructed using Python objects. SQLAlchemy automatically sanitizes inputs, making the system immune to standard SQL injection attacks.

6.3 Validation Strategies
Validation is performed at both the Client-side (Frontend) and Server-side (Backend) to ensure robust error handling and data consistency.
A. Client-Side Validation (React.js):
•	Immediate Feedback: Uses HTML5 attributes (required, type="email") and JavaScript logic to catch errors before data is sent to the server.
•	State Checks: React state management ensures that a user cannot click "Book Ticket" if the passenger count is 0 or if the source and destination are the same.
•	UX Enhancement: Loading spinners are displayed during data fetching to prevent double-submissions.
B. Server-Side Validation (Pydantic Schemas):
•	Data Integrity: The backend uses Pydantic models to strictly enforce data types.
Example Model:
Python
class TicketCreate(BaseModel):
    bus_id: str
    count: int
    source: str
    destination: str

    @validator('count')
    def check_count(cls, v):
        if v < 1:
            raise ValueError('Ticket count must be at least 1')
        return v
•	Logic: If a malicious user bypasses the frontend and sends a request with count: -5, Pydantic intercepts it and returns a 422 Unprocessable Entity error, protecting the database logic.
•	Business Logic Validation: The backend checks the database state (e.g., "Is the bus actually active?", "Is the seat capacity exceeded?") before confirming any transaction.























CHAPTER 7

Conclusion
The Bus Kar Bhai project successfully demonstrates the design and implementation of a cloud-native, centralized public transport management system that addresses key limitations of traditional fleet operations. By integrating real-time GPS tracking, dynamic route allocation, and digital ticketing into a single cohesive platform, the system improves operational efficiency, reduces passenger wait times, and ensures transparent fleet management.
The project emphasizes strong database design principles through the use of a well-structured Entity–Relationship model, a normalized relational schema up to Third Normal Form (3NF), and clearly defined functional dependencies. These design choices eliminate data redundancy, prevent update anomalies, and ensure data consistency across all operations. The use of a cloud-hosted PostgreSQL database ensures reliable transaction processing for ticketing and schedule management, even under high-concurrency scenarios.
In addition, the integration of a microservices architecture using Docker containers demonstrates the practical application of modern cloud computing concepts. The system effectively combines static schedule data with high-velocity live tracking events, showcasing a robust approach to handling real-time data streams. Role-based access control (RBAC) further ensures data security and privacy for passengers, drivers, and administrators.
Overall, Bus Kar Bhai provides a scalable, reliable, and user-centric solution for modern urban mobility. It simplifies the commute for passengers, empowers administrators with actionable insights, and demonstrates the effective application of advanced Database Management System (DBMS) and Cloud Computing concepts in a real-world transportation scenario.
.
.










     References
[1] R. Elmasri and S. B. Navathe, Fundamentals of Database Systems, 7th ed. Boston, MA,
USA: Pearson, 2016.
[2] A. Silberschatz, H. F. Korth, and S. Sudarshan, Database System Concepts, 6th ed. New
York, NY, USA: McGraw-Hill, 2011.
[3] C. J. Date, An Introduction to Database Systems, 8th ed. Boston, MA, USA: Pearson
Education, 2004.
[4] T. Connolly and C. Begg, Database Systems: A Practical Approach to Design,
Implementation, and Management, 6th ed. Harlow, U.K.: Pearson, 2015.
[5] PostgreSQL Global Development Group, “PostgreSQL Documentation,” 2023. [Online].
Available: https://www.postgresql.org/docs/
Accessed: Jan. 2026.
[6] M. Fowler, UML Distilled: A Brief Guide to the Standard Object Modeling Language,
3rd ed. Boston, MA, USA: Addison-Wesley, 2004.
[7] IEEE Computer Society, “IEEE Standard Glossary of Software Engineering
Terminology,” IEEE Std 610.12-1990, 1990.
[8] Pressman, R. S., and B. R. Maxim, Software Engineering: A Practitioner’s Approach, 8th
ed. New York, NY, USA: McGraw-Hill Education, 2015.





                                                              Appendix
Screenshots with description

 
Figure 5: Login Page 
The secure entry point for all users. It features a unified interface that toggles between "Sign In" and "Sign Up" modes. It handles role-based authentication (Passenger, Driver, Admin) using JWT tokens to ensure secure access to specific dashboards.

 
Figure 6: Admin Dashboard
The central command center for fleet management. It provides a real-time overview of the entire bus network, allowing administrators to monitor active buses, view crowded routes, and manage driver assignments. It acts as the "God's Eye View" for the system.





          
Figure 7: Live Map 
An interactive map interface powered by Leaflet.js. It visualizes the real-time GPS locations of all active buses. Passengers use this to track their bus's exact position, reducing anxiety and wait times.
 

                                                                  
Figure 8: Trip Planner 
A user-friendly search engine for passengers. Users enter their source and destination, and the system queries the database to display the best available bus routes, estimated travel times, and current seat availability.


 
Figure 9: Driver Dashboard 
A mobile-responsive interface designed for bus drivers. It displays their assigned route, current shift status, and upcoming stops. It replaces manual logbooks with a digital workflow.


 
Figure 10:  ETM Controller 
The digital ticketing module used by drivers. It allows them to issue tickets to passengers instantly. Crucially, every ticket sold updates the live "Occupancy" count of the bus, feeding real-time crowd data back to the server.
 
Figure 11: Admin Analytics 
A data visualization tool for administrators. It processes historical data to generate insights on peak traffic hours, most profitable routes, and fleet utilization, helping in long-term strategic planning.



	               






